\documentclass[10pt,conference]{IEEEtran}

\usepackage{amsmath,amssymb,amsthm}
\usepackage{graphicx}
\usepackage{subcaption}
\usepackage{booktabs}
\usepackage{xcolor}
\usepackage{cite}
\usepackage{bm}
\usepackage{hyperref}
\usepackage{algorithm}
\usepackage{algorithmic}
\usepackage{array}
\usepackage{multirow}
\usepackage{url}
\usepackage[normalem]{ulem}
\usepackage{mathtools}
\usepackage{threeparttable}

\usepackage[normalem]{ulem} % For text strikethrough


%% ── Theorem environments ──────────────────────────────────────────────────
\newtheorem{theorem}{Theorem}
\newtheorem{proposition}[theorem]{Proposition}
\newtheorem{corollary}[theorem]{Corollary}
\newtheorem{lemma}[theorem]{Lemma}
\newtheorem{definition}{Definition}
\newtheorem{assumption}{Assumption}
\newtheorem{remark}{Remark}
\newtheorem{problem}{Problem}

%% ── Macros ────────────────────────────────────────────────────────────────
\newcommand{\calR}{\mathcal{R}}
\newcommand{\calT}{\mathcal{T}}
\newcommand{\calG}{\mathcal{G}}
\newcommand{\calN}{\mathcal{N}}
\newcommand{\calC}{\mathcal{C}}
\newcommand{\calO}{\mathcal{O}}
\newcommand{\calE}{\mathcal{E}}
\newcommand{\bsig}{\bm{\sigma}}
\newcommand{\bx}{\mathbf{x}}
\newcommand{\bu}{\mathbf{u}}
\newcommand{\bp}{\mathbf{p}}
\newcommand{\bv}{\mathbf{v}}
\newcommand{\ba}{\mathbf{a}}
\newcommand{\muij}{\mu_{ij}}
\newcommand{\mutij}{\tilde{\mu}_{ij}}
\newcommand{\norm}[1]{\left\|#1\right\|}
\newcommand{\eps}{\varepsilon}

\newcommand{\ora}[1]{\textcolor{orange}{#1}}
\newcommand{\brow}[1]{\textcolor{brown}{#1}}
\newcommand{\red}[1]{\textcolor{red}{#1}}
\newcommand{\bl}[1]{\textcolor{blue}{#1}}
\newcommand{\yel}[1]{\textcolor{yellow}{#1}}
\newcommand{\shubham}[1]{\textcolor{cyan}{#1}}

%% ══════════════════════════════════════════════════════════════════════════
\begin{document}

\title{VISTA: Metric-Scale-Free Visual Control for Soft-Precise Landing under Visibility Constraints}

\maketitle
\begin{abstract}
This paper presents Visibility-Integrated Soft-Touchdown Adaptive Sliding-Mode Control (VISTA), a monocular visual-control framework for soft-precise landing on unpredictably moving platforms. Image-position and image-velocity regulation are intrinsically coupled through the monocular visual dynamics. The lateral image dynamics depend on both lateral and vertical image velocity, while regulation of the vertical component provides the mechanism for soft touchdown without metric depth. VISTA exploits this structure through a metric-scale-free architecture that couples image-position and image-velocity errors within a prescribed-performance sliding manifold. Depth-normalized relative velocity is recovered directly from optical flow, enabling simultaneous lateral alignment and vanishing relative touchdown velocity without metric target-depth, metric UAV position/velocity, or target-velocity estimation. Image-derived yaw regulation closes the orientation-alignment loop without an absolute-heading measurement. A visibility-constrained barrier conditions the commanded body lean to preserve target visibility. Inertial measurements for virtual-camera compensation and attitude realization. Lyapunov analysis establishes uniform ultimate boundedness of the coupled visual dynamics, preservation of the prescribed visual-error envelopes under the derived coupling condition, and conditional visibility invariance. Simulations achieve soft-precise touchdown for all $\boldsymbol{25}$ tested initial-condition/motion combinations and all $\boldsymbol{20}$ runs over a $\boldsymbol{\pm40\%}$ target-speed sweep, while four representative visual landing baselines do not achieve the same combined capability.
\end{abstract}

\section{INTRODUCTION}
\label{intro: section}

Autonomous landing on moving platforms requires accurate lateral alignment, low relative touchdown velocity, and retention of the target within the camera field of view (FoV)~\cite{cho2022}. Under monocular feedback, image-position and image-velocity regulation are intrinsically coupled. The lateral image dynamics depend on both lateral and vertical image motion, while attitude changes required for alignment also affect target visibility. Consequently, precise alignment, image-velocity regulation for soft touchdown, and visibility preservation cannot generally be treated as independent control objectives. The challenge is amplified when the platform state is unavailable and the landing loop must be closed from onboard visual and inertial measurements alone.

Monocular vision is attractive for size-, weight-, power-, and cost-constrained (SWaP-C) UAVs because it avoids dedicated ranging hardware. Position-based visual servoing (PBVS) typically closes the loop after metric pose reconstruction~\cite{lin2022}, whereas image-based visual servoing (IBVS) regulates image measurements directly~\cite{chaumette2006,lee2012,xie2020,lin2023}. In a monocular formulation, the image-feature and optical-flow dynamics remain coupled through depth-normalized motion. This also creates an opportunity for metric-scale-free control: optical flow provides depth-normalized relative velocity, so the inverse-depth scale need not be reconstructed explicitly and can instead be treated as uncertain control effectiveness.

Optical-flow landing methods primarily regulate touchdown motion~\cite{izzo2011,herisse2012,ho2018,singhal2025}, while other vision-based moving-platform approaches emphasize lateral alignment, relative-state tracking, or trajectory tracking~\cite{lin2022,zhao2021,zhang2026,lin2023,cho2022}. Visibility constraints have also been treated at the visual-servo level using barrier functions~\cite{salehi2021}. Landing-specific visibility-aware control has been addressed through CBF--QPs~\cite{huang2022}, as well as through predictive and perception-aware formulations~\cite{yang2025,angelis2026}. Existing visual landing approaches typically focus on either image alignment, optical-flow-based touchdown, or visibility preservation. The coupled visual regulation problem arising from simultaneous
image-position and image-velocity regulation for precise alignment and soft touchdown under visibility constraints has received comparatively little attention.

VISTA addresses this coupled problem through a metric-scale-free visual architecture that regulates depth-normalized relative motion directly. Rather than reconstructing metric target depth or target motion, the unknown inverse-depth scale is absorbed into the visual dynamics as uncertain control effectiveness. The translational and image-orientation objectives are then closed through visual observables, with inertial measurements used for virtual-camera compensation and attitude realization. The resulting capability is simultaneous lateral alignment, soft touchdown, image-derived orientation alignment, and visibility preservation without metric target-depth, metric UAV position/velocity, target-velocity, or an absolute-heading measurement.

The main contributions are:
\begin{enumerate}
    \item A coupled, metric-scale-free visual landing architecture that exploits the intrinsic coupling between image-position and image-velocity dynamics. Their joint regulation provides lateral alignment and soft touchdown without metric target-depth, metric UAV position/velocity, or target-velocity estimation.
    \item A complete image-observable landing loop that augments the coupled translational regulation with orientation alignment and visibility preservation. This closes the orientation objective without an external absolute-heading reference and preserves target visibility within the available control authority.
    \item Closed-loop guarantees establishing uniform ultimate boundedness of the coupled visual sliding dynamics and image-orientation error, preservation of the prescribed visual-error envelopes under the derived coupling condition, and conditional forward invariance of the defined visibility set. Simulations validate the resulting capability across multiple initial conditions, target motions, speed variations, and four visual landing baselines.
\end{enumerate}

Table~\ref{tab:comparison} summarizes the capabilities relevant to the coupled landing problem for these baselines.

\begin{table}[!t]
\centering
\caption{Comparison with the four visual landing controllers used as
baselines.}
\label{tab:comparison}
\scriptsize
\setlength{\tabcolsep}{4pt}
\begin{threeparttable}
\begin{tabular}{lcccc}
\toprule
\textbf{Method} &
\textbf{Soft} &
\textbf{Visibility} &
\textbf{Metric-Depth} &
\textbf{Camera+IMU} \\
&
\textbf{Touchdown} &
\textbf{Constraint} &
\textbf{Free} &
\textbf{Only} \\
\midrule
PBVS--PPC~\cite{lin2022}
    & $\times$   & $\times$   & $\times$   & $\times$   \\
PBVS--AEDO~\cite{zhang2026}
    & $\times$   & $\times$   & $\times$   & $\times$   \\
IBVS--PPC~\cite{lin2023}
    & $\times$   & $\times$   & \checkmark & \checkmark \\
FF--IBVS~\cite{cho2022}
    & $\times$   & $\circ$    & $\times$   & $\times$   \\
\textbf{VISTA}
    & \checkmark & \checkmark & \checkmark & \checkmark \\
\bottomrule
\end{tabular}
\begin{tablenotes}[flushleft]
\footnotesize
\item \checkmark: explicit provision; $\times$: not established; $\circ$: partial/heuristic treatment. ``Metric-Depth Free'' means no online metric target-depth estimation is required; ``Camera+IMU Only'' denotes operation without external localization or target-state sensing. The VISTA visibility entry denotes explicit visibility-constrained control.
\end{tablenotes}
\end{threeparttable}
\end{table}


\section{PROBLEM FORMULATION}
\label{sec:problem}
Consider a quadrotor equipped with a downward-facing monocular camera and an inertial measurement unit (IMU), tasked with landing on a moving planar target. The controller has access only to the image and inertial measurements, while the metric target depth, absolute yaw angle, metric UAV position/velocity, and the translational and rotational target states are unavailable.

\subsection{Quadrotor and Visual Model}
\label{subsec:dynamics_model}
Let $\mathcal{I}$ denote the inertial frame and $\mathcal{B}=\{O_{\text{b}},X_{\text{b}},Y_{\text{b}},Z_{\text{b}}\}$ denote the body-fixed frame, with linear and angular velocities ${}^\mathcal{B}\boldsymbol{v}_\text{b}$ and ${}^\mathcal{B}\boldsymbol{\omega}_\text{b}$. From the Newton--Euler formulation, the flight dynamics of the quadrotor are
\begin{equation}
\label{quadrotor dynamics: equation}
    \begin{aligned}
    {}^\mathcal{B}\dot{\boldsymbol{v}}_\text{b} &= \tfrac{1}{m}\big(-{}^\mathcal{B} T_u \hat{\boldsymbol{e}}_z + {}^\mathcal{B}\boldsymbol{F}_\text{g} + {}^\mathcal{B}\boldsymbol{F}_\text{d}\big) - {}^\mathcal{B}\boldsymbol{\omega}_\text{b}\times{}^\mathcal{B}\boldsymbol{v}_\text{b},\\
    J {}^\mathcal{B}\dot{\boldsymbol{\omega}}_\text{b} &= {}^\mathcal{B}\boldsymbol{\tau}_u + {}^\mathcal{B}\boldsymbol{\tau}_\text{d} - {}^\mathcal{B}\boldsymbol{\omega}_\text{b}\times J {}^\mathcal{B}\boldsymbol{\omega}_\text{b},
    \end{aligned}
 \end{equation}
where $m$ and $J$ denote the quadrotor mass and inertia matrix, respectively, $\hat{\boldsymbol{e}}_z=[0,0,1]^\top$ is the unit vector  along $+Z_{\mathrm b}$, $\boldsymbol{u}=[{}^\mathcal{B}T_u,{}^\mathcal{B}\boldsymbol{\tau}_u^\top]^\top\in\mathbb{R}^4$ is the combined thrust--torque input, with ${}^\mathcal{B}T_u \geq 0$ the rotor-thrust magnitude acting along $-Z_{\mathrm b}$, ${}^\mathcal{B} \boldsymbol{F}_\text{g}$ is the gravitational force, and ${}^\mathcal{B} \boldsymbol{F}_\text{d}$ and ${}^\mathcal{B} \boldsymbol{\tau}_\text{d}$ collect unknown bounded disturbances.

Let $\mathcal C$ denote the downward-facing camera frame. To remove the apparent feature motion induced by the UAV roll and pitch, the measured image feature points are reprojected onto a virtual camera frame $\mathcal V$, constructed by compensating the UAV roll and pitch such that its image plane remains parallel to the inertial horizontal plane, as illustrated in Fig.~\ref{frames planes: figure}.

\begin{figure}[!t]
\centering
    \includegraphics[width=0.8\columnwidth]{Figures/frames_planes.pdf}
    \caption{\textbf{Coordinate frames and virtual-image geometry.} The roll--pitch-compensated virtual frame $\mathcal V$ has an image plane parallel to the inertial $X_{\mathrm i}$--$Y_{\mathrm i}$ plane; $O_{\mathrm b}=O_{\mathrm c}=O_{\mathrm v}$ in the model (offsets shown only for clarity). A target point and its camera/virtual projections are shown.}
    \label{frames planes: figure}
\end{figure}

\begin{assumption}[Feature configuration]
\label{asm:feature_config}
The $N\geq3$ feature points are rigidly attached to the target and lie non-collinearly on a common plane. The feature centroid coincides with the target rotational reference point, and the camera optical center coincides with the UAV body origin, so that rotational lever-arm velocities do not enter the centroid kinematics.
\end{assumption}

\begin{assumption}[Nominal depth geometry]
\label{asm:depth_geometry}
For the nominal level-platform geometry, the target plane is parallel to the virtual image plane, so that the reprojected feature points share a common virtual-frame depth $z(t)>0$. Bounded platform roll and pitch introduce bounded unequal-depth perturbations, whose induced visual modeling errors are included in the lumped disturbances.
\end{assumption}

Let ${}^\mathcal{C}\hat{\boldsymbol{r}}_i$ denote the feature coordinates on the camera image plane and $\hat{\boldsymbol{r}}_i=[\hat{x}_i,\hat{y}_i]^\top, i\in\{1,\ldots,N\}$ denote their reprojection onto the virtual image plane with $\hat{\boldsymbol{r}} = \frac{1}{N}\sum_{i=1}^N\hat{\boldsymbol{r}}_i$ their centroid.
The centroid is used as the image-position variable for lateral alignment.

\subsection{Monocular Visual Observables}
Let $\boldsymbol{r}_\text{t}\in\mathbb R^3$ denote the position of the target feature centroid relative to the UAV, expressed in $\mathcal V$. 
Under Assumption~\ref{asm:depth_geometry}, the focal-length-normalized centroid defines the homogeneous image-position variable
\begin{equation}
    \boldsymbol{s} = \begin{bmatrix}\hat{\boldsymbol{r}}/f\\ 1\end{bmatrix} = \lambda(t)\,\boldsymbol{r}_\text{t}, \qquad \lambda(t)\triangleq\frac{1}{z(t)}>0,
    \label{image position: equation}
\end{equation}
where $z(t)$ is the camera--target depth.
Define the translational and angular image velocities as
\begin{equation}
    \boldsymbol{\nu}\triangleq\lambda(t)\boldsymbol v_{\mathrm{rel}}, \qquad
    \boldsymbol w \triangleq \boldsymbol\omega_{\mathrm t} - \dot\psi_{\mathrm b}\hat{\boldsymbol e}_z,
    \label{image velocities: equation}
\end{equation}
where $\boldsymbol v_{\mathrm{rel}}\in\mathbb R^3$ denotes the target translational velocity relative to the UAV expressed in $\mathcal V$, $\boldsymbol\omega_{\mathrm t}$ is the unknown target angular velocity
and $\dot\psi_{\mathrm b}$ is the UAV yaw rate.

For the nominal level-platform geometry of Assumption~\ref{asm:depth_geometry}, the optical flow of each virtual-image feature satisfies $\dot{\hat{\boldsymbol{r}}}_i = L_{s,i}^\text{t}\boldsymbol{\nu} + L_{s,i}^\text{r}\boldsymbol{w}$, where
\begin{equation*}
L_{s,i}^\text{t} = \begin{bmatrix} f & 0 & -\hat{x}_i \\ 0 & f & -\hat{y}_i \end{bmatrix}, \quad
L_{s,i}^\text{r} = \begin{bmatrix} -\frac{\hat{x}_i\hat{y}_i}{f} & \frac{f^2+\hat{x}_i^2}{f} & -\hat{y}_i \\ -\frac{f^2+\hat{y}_i^2}{f} & \frac{\hat{x}_i\hat{y}_i}{f} & \hat{x}_i \end{bmatrix}.
\end{equation*}
Since the inverse depth is absorbed into $\boldsymbol\nu$, $L_{s,i}^\text{t}$ and $L_{s,i}^\text{r}$ depend only on the measured image coordinates and known focal length. Hence, $\boldsymbol\nu$ and $\boldsymbol w$ are recovered from optical flow without metric-depth reconstruction. Defining $L_{s,i}=[L_{s,i}^\text{t}~~L_{s,i}^\text{r}]$ and stacking the $N$ feature equations gives $L_s=[L_{s,1}^\top \cdots L_{s,N}^\top]^\top\in\mathbb{R}^{2N\times 6}$.

\begin{assumption}[Optical-flow observability]
\label{asm:flow_observability}
The stacked interaction matrix $L_s$ has full column rank
over the considered feature configurations.
\end{assumption}

Under Assumption~\ref{asm:flow_observability},
\begin{equation}
    \begin{bmatrix} \boldsymbol{\nu}^\top~\boldsymbol{w}^\top\end{bmatrix}^\top = L_s^\dagger\begin{bmatrix}\dot{\hat{\boldsymbol{r}}}_1^\top \cdots \dot{\hat{\boldsymbol{r}}}_N^\top\end{bmatrix}^\top.
    \label{optic flow inverse: equation}
\end{equation}
The virtual-image target orientation is defined by the principal-axis $\alpha =  \frac{1}{2}\operatorname{atan2}\big(2\mu_{11},\mu_{20}-\mu_{02}\big)$, where $\mu_{pq} = \sum_{i=1}^{N}(\hat{x}_i - \hat{x})^p(\hat{y}_i - \hat{y})^q$ are the centered image moments.

\begin{assumption}[Image-orientation observability]
\label{asm:orientation_observability}
During the landing interval, the target second-moment tensor remains anisotropic such that $\left(\mu_{20}-\mu_{02}\right)^2 + 4\mu_{11}^2 \ge c_\alpha, ~ c_\alpha>0$, which excludes rotationally symmetric feature configurations.
\end{assumption}

Using the virtual-camera kinematics under Assumptions~\ref{asm:feature_config},
\ref{asm:depth_geometry}, and \ref{asm:orientation_observability}, the nominal visual dynamics are:
\begin{equation} \label{visual dynamics: equation}
    \begin{aligned}
    \dot{\boldsymbol{s}} &= \boldsymbol{\nu} - \dot{\psi}_\text{b}\,\hat{\boldsymbol{e}}_z\times\boldsymbol{s} - (\boldsymbol{\nu}\cdot\hat{\boldsymbol{e}}_z)\,\boldsymbol{s}, \\
    \dot{\boldsymbol{\nu}} &= \lambda(t)\boldsymbol{a} - \dot{\psi}_\text{b}\,\hat{\boldsymbol{e}}_z\times\boldsymbol{\nu} - (\boldsymbol{\nu}\cdot\hat{\boldsymbol{e}}_z)\,\boldsymbol{\nu}, \\
    \dot{\alpha} &= \boldsymbol l_\alpha^\top\boldsymbol w, \qquad l_{\alpha,3}=1.
    \end{aligned}
\end{equation}
Here $\boldsymbol{a}$ denotes the relative translational acceleration expressed in $\mathcal{V}$, and $\boldsymbol l_\alpha$ is the orientation interaction vector obtained by differentiating the moment-based principal-axis angle with respect to the feature motion. Its components are algebraic functions of the centered image moments and satisfy $l_{\alpha,3}=1$.

\begin{assumption}[Visual-signal regularity]
\label{asm:visual_regularity}
Over the landing interval, the visual observables $\boldsymbol s$, $\alpha$, $\boldsymbol\nu$, and $\boldsymbol w$, together with the time derivatives required by the controller and stability analysis, remain bounded.
\end{assumption}

\subsection{Problem Statement}
\label{subsec:problem_statement}
The relative translational acceleration is modeled as $\boldsymbol{a} = -\beta_{\mathrm{a}}(t)\boldsymbol{a}_{\mathrm{u}} + \boldsymbol{d}_{\mathrm{a}}$, where $\boldsymbol{d}_{\mathrm{a}}\in\mathbb{R}^{3}$ collects the unknown target acceleration and additive translational disturbances. The negative sign follows from defining the target motion relative to the UAV.

Since the inverse-depth scale $\lambda(t)$ is unknown, define the effective input gain $\beta(t) \triangleq\lambda(t)\beta_{\mathrm{a}}(t)>0$. Then,
\begin{equation}
    \dot{\boldsymbol{\nu}} = -\beta(t)\boldsymbol{a}_{\mathrm{u}} -\dot{\psi}_{\mathrm{b}}
    \hat{\boldsymbol{e}}_{z}\times\boldsymbol{\nu} -(\boldsymbol{\nu}\cdot\hat{\boldsymbol{e}}_{z})\boldsymbol{\nu} + \boldsymbol{d}_{\nu},
    \label{uncertain image velocity dynamics: equation}
\end{equation}
where $\boldsymbol d_\nu$ denotes the lumped image-space disturbance, including $\lambda(t)\boldsymbol d_a$ and the bounded visual-modeling perturbations associated with the unequal-depth effects of Assumption~\ref{asm:depth_geometry}.
The corresponding image-orientation dynamics are
\begin{equation}
    \dot{\alpha} = -\omega_{\mathrm{u}} + d_{\alpha}, \qquad
    d_{\alpha} \triangleq \boldsymbol l_\alpha^\top \boldsymbol\omega_{\mathrm t} + (\omega_{\mathrm{u}} - \dot{\psi}_{\mathrm{b}}),
    \label{uncertain image orientation dynamics: equation}
\end{equation}
where $\omega_{\mathrm{u}}$ is the commanded yaw-rate, and $d_{\alpha}$ collects the target-induced image-orientation motion and yaw-rate realization error.

\begin{assumption}[Bounded uncertainties]
\label{asm:uncertainty}
During the landing interval, the camera-target depth $z(t)$ remains strictly positive, and the effective input gain satisfies $0<\underline{\beta} \leq \beta(t) \leq \overline{\beta}$, where $\underline{\beta}$ and $\overline{\beta}$ are finite but unknown constants. Moreover, the lumped additive uncertainties satisfy $\|\boldsymbol{d}_{\nu}(t)\| \leq \bar{d}_{\nu}, ~ |d_{\alpha}(t)| \leq \bar{d}_{\alpha}$, for finite but unknown constants $\bar{d}_{\nu}>0$ and $\bar{d}_{\alpha}>0$. In addition, $d_{\alpha}$ is continuously differentiable with bounded derivative.
\end{assumption}

Define $\Phi \triangleq \operatorname{diag} \left( \varphi_{x,\max}, \varphi_{y,\max} \right)$, where $\varphi_{x,\max}$ and $\varphi_{y,\max}$ denote the focal-length-normalized FoV half-widths along the two image axes.
Since $s_z\equiv1$, define the FoV-normalized image-position error as $\boldsymbol{e}_{p} \triangleq \Phi^{-1} \left( \boldsymbol{s}_{xy} - \boldsymbol{s}_{\mathrm{d},xy} \right) \in\mathbb{R}^{2}$, where $\boldsymbol{s}_{\mathrm{d},xy}$ is the desired normalized image location of the target centroid.
Define the image-orientation error as $e_{\alpha} = \frac{1}{2} \operatorname{atan2} \left(\sin\!\left(2(\alpha-\alpha_{\mathrm{d}})\right),\cos\!\left(2(\alpha-\alpha_{\mathrm{d}})\right)\right)\in\left(-\frac{\pi}{2},\frac{\pi}{2}\right]$, where $\alpha_{\mathrm{d}}$ is the desired target orientation in the virtual image plane.

Target visibility imposes an additional physical constraint that is not guaranteed by image-position regulation. Let ${}^{\mathcal C}\tilde{\boldsymbol r} = [\tilde{x},\tilde{y}]^\top$ denote the measured marker centre on the camera image plane in focal-length-normalized coordinates. Then, the visibility set is 
\begin{equation}
    \mathcal S_{\mathrm{vis}} = \left\{{}^{\mathcal C}\tilde{\boldsymbol r}\in\mathbb R^2:\left\| \Phi^{-1}{}^{\mathcal C}\tilde{\boldsymbol r} \right\|_\infty \le 1 \right\}.
    \label{visibility set: equation}
\end{equation}

Since $\boldsymbol v_{\mathrm{rel}}=z\boldsymbol\nu$, bounded regulation of $\boldsymbol\nu$ makes the metric relative velocity vanish with depth, providing the soft-touchdown mechanism without metric-depth or metric relative-velocity estimation. A physical interpretation follows from the relative translational kinetic energy. Let $m_{\mathrm u}$ and $m_{\mathrm p}$ denote the UAV and platform masses and define the reduced mass $\mu=m_{\mathrm u}m_{\mathrm p}/(m_{\mathrm u}+m_{\mathrm p})$. Then
\begin{equation}
E_{\mathrm{rel}}=\frac12\mu\|\boldsymbol v_{\mathrm{rel}}\|^2
=\frac12\mu z^2\|\boldsymbol\nu\|^2,
\label{relative touchdown energy: equation}
\end{equation}
so $E_{\mathrm{rel}}=\mathcal O(z^2)\to0$ as $z\to0$ for bounded $\boldsymbol\nu$. Thus, the image-space velocity regulation implies both vanishing relative touchdown velocity and vanishing relative-motion energy.

\begin{problem}
\label{prob:landing}
Using only monocular image measurements and the available inertial measurements, design a visual feedback controller such that the image-position and image-velocity errors remain within prescribed performance envelopes, the image-orientation error is ultimately bounded, regulation of the vertical image velocity toward the reference $\nu_{\mathrm d,z}<0$ yields vanishing relative touchdown velocity, and the visibility set $\mathcal S_{\mathrm{vis}}$ is forward invariant, without requiring metric target depth, metric UAV position/velocity, absolute yaw, or target-state measurements.
\end{problem}

\section{COUPLED IMAGE-FEATURE AND OPTICAL-FLOW CONTROL}
\label{sec:control}
The controller is constructed directly in the visual domain, following the intrinsic coupling of the visual dynamics rather than treating image-position and image-velocity regulation as independent loops. A prescribed-performance transformation constrains the normalized image-position error, from which a compatible image-velocity reference is generated. The transformed image-position and image-velocity errors are then coupled through a common sliding manifold so that lateral alignment and relative-velocity convergence are enforced coherently in metric-scale-free coordinates (Fig.~\ref{fig:architecture}).

For any constrained vector $\boldsymbol q\in\mathbb R^n$, consider the following component-wise prescribed-performance transformation: 
\begin{equation}
    \begin{aligned}
    -\rho_{q,k}(t)&<q_k(t)<\rho_{q,k}(t), \qquad k\in\{1,\ldots,n\} \\ 
    \rho_{q,k}(t) &= (\rho_{q,k,0}-\rho_{q,k,\infty}) e^{-\ell_{q,k}t} + \rho_{q,k,\infty},  \quad \ell_{q,k}>0\\
    q_k &= \rho_{q,k}(t) S_{q,k}, \qquad
    S_{q,k} = \tanh\!\left(\frac{\zeta_{q,k}}{2}\right).
    \end{aligned}
    \label{generic PPC transformation: equation}
\end{equation}
For any initial condition satisfying $|q_k(0)|<\rho_{q,k}(0)$, the transformation maps the constrained error $q_k$ to the unconstrained variable $\zeta_{q,k}\in\mathbb R$. Differentiating \eqref{generic PPC transformation: equation} gives
\begin{equation}
    \begin{aligned}
    \dot{\boldsymbol\zeta}_q &= \mathcal G_q \left[ \dot{\boldsymbol q} - \mathcal S_q\dot{\boldsymbol\rho}_q \right], \quad \mathcal S_q = \operatorname{diag}(S_{q,1},\ldots,S_{q,n}), \\
    \mathcal G_q &= \operatorname{diag}(g_{q,1},\ldots,g_{q,n}), \quad
    g_{q,k} = \frac{(e^{\zeta_{q,k}}+1)^2} {2e^{\zeta_{q,k}}\rho_{q,k}}.
    \end{aligned}
    \label{generic transformed error dynamics: equation}
\end{equation}

\subsection{Coupled Translational Visual Control}
\label{subsec:translational_control}
From \eqref{visual dynamics: equation}, for a constant desired image location $\boldsymbol s_{\mathrm{d},xy}$, the normalized image-position error satisfies
\begin{equation}
    \dot{\boldsymbol e}_{p} = \Phi^{-1} \left[\boldsymbol\nu_{xy} - \left(\dot{\psi}_{\mathrm b}\hat{\boldsymbol e}_{z}\times\boldsymbol s\right)_{xy} - \nu_{z}\boldsymbol s_{xy} \right].
    \label{image position error dynamics: equation}
\end{equation}
Applying \eqref{generic PPC transformation: equation} to $\boldsymbol e_p$ yields
\begin{equation}
    \dot{\boldsymbol\zeta}_{p} = \mathcal G_{p} \left[ \dot{\boldsymbol e}_{p} - \mathcal S_{p}\dot{\boldsymbol\rho}_{p}(t) \right],
    \label{transformed image error dynamics: equation}
\end{equation}
so that $\rho_{p,k}$ specifies the admissible fraction of the available image half-width along each image axis.

Choose the compatible translational image-velocity reference as
\begin{equation}
    \boldsymbol\nu_{\mathrm d} = \begin{bmatrix} \Phi\left(\mathcal S_p\dot{\boldsymbol\rho}_p - k_p\mathcal G_p^{-1}\boldsymbol\zeta_p \right)\\[1mm] 0 \end{bmatrix} + \dot{\psi}_{\mathrm b} \hat{\boldsymbol e}_{z}\times\boldsymbol s + \nu_{\mathrm d,z}\boldsymbol s,
    \label{desired image velocity: equation}
\end{equation}
where $\nu_{\mathrm d,z}<0$ is the prescribed vertical image-velocity reference that governs the touchdown approach, and $k_p>0$ is the image-position convergence gain.
Since $\nu_z=v_{\mathrm{rel},z}/z=\dot z/z$, ideal regulation to the constant reference $\nu_{\mathrm d,z}<0$ yields $z(t)=z(0)e^{\nu_{\mathrm d,z}t}$ and hence $v_{\mathrm{rel},z}=\nu_{\mathrm d,z}z(t)\to0$.
Moreover, because $\boldsymbol v_{\mathrm{rel}}=z\boldsymbol\nu$, bounded regulation of the complete image velocity $\boldsymbol\nu$ causes all components of the target-relative velocity to vanish as $z\to0$, providing the soft-touchdown mechanism without metric-depth estimation.

Define the image-velocity error as $\boldsymbol e_{\nu} \triangleq \boldsymbol\nu-\boldsymbol\nu_{\mathrm d}$.
Substitution into \eqref{transformed image error dynamics: equation} yields
\begin{equation}
    \dot{\boldsymbol\zeta}_p = - k_p\boldsymbol\zeta_p + \mathcal G_p\Phi^{-1} \left( \boldsymbol e_{\nu,xy} - e_{\nu,z}\boldsymbol s_{xy} \right).
    \label{coupled transformed position dynamics: equation}
\end{equation}
Equation~\eqref{coupled transformed position dynamics: equation} makes the central coupling explicit. The image-position dynamics depend directly on the image-velocity error, including the vertical term $e_{\nu,z}\boldsymbol{s}_{xy}$. Therefore, regulating the vertical image velocity to obtain soft touchdown also affects lateral image-position convergence whenever the target is
off-center. Image-position and image-velocity regulation cannot, in general, be designed as independent channels. The compatible reference \eqref{desired image velocity: equation} and coupled manifold \eqref{coupled sliding manifold: equation} are introduced to regulate this interaction directly.
Under Assumption~\ref{asm:visual_regularity} and the prescribed smooth performance functions, $\boldsymbol\nu_{\mathrm d}\in C^1$ and $\dot{\boldsymbol\nu}_{\mathrm d}$ remains bounded over the landing interval.

Using \eqref{uncertain image velocity dynamics: equation}, the translational image-velocity error satisfies
\begin{equation}
    \dot{\boldsymbol e}_{\nu} = -\beta(t)\boldsymbol a_{\mathrm u} + \boldsymbol c_{\nu} + \boldsymbol d_{\nu},
    \label{image velocity error dynamics: equation}
\end{equation}
where $\boldsymbol c_{\nu} \triangleq -\dot{\psi}_{\mathrm b}\hat{\boldsymbol e}_{z}\times\boldsymbol\nu -\nu_z\boldsymbol\nu -\dot{\boldsymbol\nu}_{\mathrm d}$.
Applying \eqref{generic PPC transformation: equation} to $\boldsymbol e_{\nu}$ yields
\begin{equation}
    \dot{\boldsymbol\zeta}_{\nu} = \mathcal G_{\nu} \left[ -\beta(t)\boldsymbol a_{\mathrm u} + \boldsymbol c_{\nu} + \boldsymbol d_{\nu} - \mathcal S_{\nu}\dot{\boldsymbol\rho}_{\nu} \right],
    \label{transformed image velocity dynamics: equation}
\end{equation}
where
$\boldsymbol\rho_\nu$, $\mathcal S_\nu$, and $\mathcal G_\nu$ follow \eqref{generic transformed error dynamics: equation} with $n=3$.

\subsection{Adaptive Control for the Coupled Sliding Manifold}
\label{subsec:coupled_surface}
Since $s_z\equiv1$, there is no independent vertical image-position error. Hence, use $\boldsymbol\zeta_p$ in the image-plane channels and the integral of $\zeta_{\nu,z}$ in the vertical channel.
Define $\boldsymbol\zeta_{\mathrm a} \triangleq \begin{bmatrix} \boldsymbol\zeta_p^\top & \int_0^t \zeta_{\nu,z}(\tau)\,d\tau \end{bmatrix}^\top \in\mathbb R^3$ for the coupled sliding manifold
\begin{equation}
    \boldsymbol\sigma = \boldsymbol\zeta_{\nu} + \boldsymbol{\chi} \,\boldsymbol\zeta_{\mathrm a}, \qquad \boldsymbol{\chi}=\operatorname{diag}(\chi_x,\chi_y,\chi_z)\succ0
    \label{coupled sliding manifold: equation}
\end{equation}
Define $\boldsymbol\eta_\nu \triangleq \boldsymbol c_\nu - \mathcal S_\nu\dot{\boldsymbol\rho}_\nu + \mathcal G_\nu^{-1}\boldsymbol{\chi}\dot{\boldsymbol\zeta}_{\mathrm a}$, where 
\begin{equation*}
\dot{\boldsymbol\zeta}_{\mathrm a} = \begin{bmatrix} - k_p\boldsymbol\zeta_p + \mathcal G_p\Phi^{-1} \left( \boldsymbol e_{\nu,xy} - e_{\nu,z}\boldsymbol s_{xy} \right) \\ \zeta_{\nu,z}
\end{bmatrix}.
\end{equation*}
Then, the coupled sliding dynamics can be written as
\begin{equation}
    \dot{\boldsymbol\sigma} = \mathcal G_\nu \left( -\beta(t)\boldsymbol a_{\mathrm u} + \boldsymbol\eta_\nu + \boldsymbol d_\nu \right).
    \label{compact sliding dynamics: equation}
\end{equation}

Define the uncertainty regressor $\mathcal R_\nu \triangleq \begin{bmatrix} - \boldsymbol\eta_\nu & d_{\nu,0} I_3 \end{bmatrix} \in\mathbb R^{3\times4}$ and the lumped uncertainty vector $\boldsymbol\delta_\nu \triangleq \underline{\beta}^{-1} \begin{bmatrix} \left(\frac{\beta(t)}{\beta_0} - 1\right) & d_{\nu,0}^{-1} \boldsymbol d_\nu^\top \end{bmatrix}^\top\in\mathbb R^4$, where $\beta_0 = 1~\mathrm{m}^{-1}$ and $d_{\nu,0}=1~\mathrm{s}^{-2}$ are fixed reference scales used to normalize the gain- and disturbance-dependent terms.
Then, the sliding dynamics can be written as
\begin{equation}
    \dot{\boldsymbol\sigma} = -\beta(t)\mathcal G_\nu\left(\boldsymbol a_{\mathrm u} - \frac{\boldsymbol\eta_\nu}{\beta_0} \right) + \underline{\beta}\,\mathcal G_\nu\mathcal R_\nu\boldsymbol\delta_\nu.
    \label{sliding uncertainty dynamics: equation}
\end{equation}
By Assumption~\ref{asm:uncertainty}, $\boldsymbol\delta_\nu$ is uniformly bounded; hence, there exists a finite but unknown constant $\tilde d_\nu>0$ such that $\|\boldsymbol\delta_\nu(t)\|_2\leq\tilde d_\nu$.
For each axis, let $\mathcal R_{\nu,k}$ denote the $k$-th row of $\mathcal R_\nu$ and define $\mathcal W_\nu \triangleq \operatorname{diag} (w_{\nu,1},w_{\nu,2},w_{\nu,3})$, where $w_{\nu,k} \triangleq \left\| \mathcal R_{\nu,k} \right\|_2$.
Hence, $\left| (\mathcal R_\nu\boldsymbol\delta_\nu)_k \right| \leq w_{\nu,k}\tilde d_\nu, k\in\{1,2,3\}$.
Since $\mathcal G_\nu$ is positive diagonal, $\left| (\mathcal G_\nu\mathcal R_\nu\boldsymbol\delta_\nu)_k \right| \le g_{\nu,k}w_{\nu,k}\tilde d_\nu$.

The commanded translational acceleration is chosen as
\begin{equation}
    \boldsymbol a_{\mathrm u} = \frac{\boldsymbol\eta_\nu}{\beta_0} + \mathcal G_\nu^{-1} \left[ \Gamma\boldsymbol\sigma + \mathcal G_\nu\mathcal W_\nu \left( \boldsymbol\kappa\odot \operatorname{sgn} (\boldsymbol\sigma) \right) \right],
    \label{adaptive acceleration control: equation}
\end{equation}
where $\Gamma\succ0$ is diagonal, $\odot$ denotes the Hadamard product, and $\operatorname{sgn}(\cdot)$ acts component-wise. At the switching surface, the ideal sign function is interpreted in the standard set-valued sliding-mode sense, $\operatorname{sgn}(0)\in[-1,1]$.
The first term compensates the known transformed visual dynamics using the nominal input scale $\beta_0$, whereas the second term provides linear stabilization and adaptive rejection of the unknown lumped uncertainty.
The adaptive switching gains $\boldsymbol\kappa = [\kappa_x,\kappa_y,\kappa_z]^\top$ evolve according to
\begin{equation}
    \dot{\boldsymbol\kappa} = \mathcal N \mathcal G_\nu \mathcal W_\nu |\boldsymbol\sigma| - \mathcal N\mathcal P\boldsymbol\kappa, \qquad \boldsymbol\kappa(0)\in\mathbb R_{>0}^{3},
    \label{adaptive gain law: equation}
\end{equation}
where $\mathcal N,\mathcal P\succ0$ are diagonal and $|\boldsymbol\sigma|$ denotes the component-wise absolute value. 
Substitution into \eqref{sliding uncertainty dynamics: equation} gives
\begin{equation}
    \dot{\boldsymbol\sigma} = -\beta(t) \left[ \Gamma\boldsymbol\sigma + \mathcal G_\nu \mathcal W_\nu\left( \boldsymbol\kappa\odot \operatorname{sgn} (\boldsymbol\sigma) \right) \right] + \underline{\beta}\mathcal G_\nu\mathcal R_\nu\boldsymbol\delta_\nu.
    \label{closed loop sliding dynamics: equation}
\end{equation}


\begin{figure*}[t]
    \centering
    \includegraphics[scale=0.4]{Drafts/ICRA/Arch_1609_6.png}
    \caption{\textbf{Closed-loop VISTA architecture and representative landing response.}
    Image and inertial sensing provide $\boldsymbol s,\alpha,\boldsymbol\nu, \boldsymbol w$. The coupled controller generates $\boldsymbol a_u$, the visibility CBF conditions its lean, and the virtual-compass loop generates $\psi_d$ for the geometric $SO(3)$ controller. The right panel shows a representative UAV--target trajectory together with the target-relative position norm, indicating landing precision, and the relative-velocity norm, indicating soft-touchdown performance.}
    \label{fig:architecture}
\end{figure*}

% \begin{figure}[t]
%     \centering
%     \includegraphics[width=\columnwidth]{Drafts/ICRA/Arch_temp3.png}
%     \caption{Architecture of the proposed method.}
%     \label{fig:architecture}
% \end{figure}

\subsection{Visibility Control Barrier Function}
\label{subsec:visibility}
The prescribed visual-error envelopes do not independently guarantee that the measured landing marker remains inside the physical camera FoV. We therefore impose a separate discrete-time visibility barrier on the measured camera-frame marker centre ${}^{\mathcal C}\tilde{\boldsymbol r}$.
For each $k\in\{1,2\}$, define the barrier function
\begin{equation}
    h_k({}^{\mathcal C}\tilde{\boldsymbol r}) = 1 - \left| \left[ \Phi^{-1} {}^{\mathcal C}\tilde{\boldsymbol r} \right]_k \right|,
    \label{cbf barrier: equation}
\end{equation}
so that $h_k\ge0$ for $k\in\{1,2\}$ is equivalent to ${}^{\mathcal C}\tilde{\boldsymbol r}\in\mathcal S_{\mathrm{vis}}$ defined in \eqref{visibility set: equation}.

Let $\boldsymbol y\in\mathbb R^2$ denote the body-lean coordinate expressed in the image axes, with nominal value $\boldsymbol y_{\mathrm d}\triangleq \boldsymbol a_{\mathrm u,xy}/|a_{\mathrm u,z}|$ and realized value $\boldsymbol y_{\mathrm{curr}}$.
Let $L_e({}^{\mathcal C}\tilde{\boldsymbol r})$ denote the local lean-to-image sensitivity and $\boldsymbol d\in\mathbb R^2$ the marker-centre drift not induced by the commanded lean, estimated directly from the de-rotated optical flow provided by the monocular front end; no target-state estimate is required.
Then, the nominal marker-centre prediction over a horizon $\tau>0$ is
\begin{equation}
    \boldsymbol c_{\mathrm{next}}(\boldsymbol y) = {}^{\mathcal C}\tilde{\boldsymbol r} + L_e({}^{\mathcal C}\tilde{\boldsymbol r}) \left(\boldsymbol y-\boldsymbol y_{\mathrm curr}\right) + \tau\boldsymbol d.
    \label{predicted marker centre: equation}
\end{equation}

The visibility-conditioned command is obtained from
\begin{equation}
    \begin{aligned}
    \boldsymbol y^\star = \arg\min_{\boldsymbol y\in\mathbb R^2} ~& \frac{1}{2}\|\boldsymbol y-\boldsymbol y_{\mathrm d}\|^2 \\
    \text{s.t.}\quad& \left| \left[ \Phi^{-1}\boldsymbol c_{\mathrm{next}}(\boldsymbol y) \right]_k \right| \le 1,\quad k\in\{1,2\},\\ 
    &\|\boldsymbol y\|\le y_{\max},
    \end{aligned}
    \label{cbf qp: equation}
\end{equation}
where $y_{\max}$ is the maximum deliverable lean-coordinate magnitude allowed by the available thrust authority.
The optimized $\boldsymbol y^\star$ replaces $\boldsymbol y_{\mathrm d}$ in the desired-attitude construction, while the nominal vertical command is retained.

\begin{assumption}[Visibility feasibility]
\label{asm:visibility_feasibility}
The visibility-constrained optimization \eqref{cbf qp: equation} remains feasible over the landing interval.
\end{assumption}

\begin{remark}[Effect of Visibility Conditioning on Translational Stability]
\label{rem:visibility_stability}
The visibility filter modifies only the nominal lateral command while retaining the vertical image-velocity command $\nu_{\mathrm d,z}$ that provides the soft-touchdown behavior. Under the finite control-authority constraint in \eqref{cbf qp: equation}, the difference between the visibility-conditioned and nominal translational accelerations is bounded and enters the image-velocity dynamics through the same matched input channel as $\boldsymbol d_\nu$. It can therefore be absorbed into the unknown disturbance bound, so Theorem~\ref{thm:translational} remains applicable to the visibility-conditioned closed loop.
\end{remark}

\subsection{Virtual-Compass Yaw and Geometric Attitude Control}
\label{subsec:orientation_control}
The image-orientation loop acts as a virtual compass, generating the desired heading without requiring an external absolute-heading reference.

For a constant desired image orientation $\alpha_{\mathrm d}$, the wrapped error satisfies $\dot e_\alpha=\dot\alpha$ away from the wrapping discontinuity. Using \eqref{uncertain image orientation dynamics: equation}, the image-orientation error dynamics are written in the lumped form
\begin{equation}
    \dot e_\alpha = -\omega_{\mathrm{u}} + d_\alpha .
    \label{image orientation error dynamics: equation}
\end{equation}
The commanded yaw rate is generated by
\begin{equation}
    \dot\omega_{\mathrm{u}} = k_\alpha e_\alpha + k_w w_z,
    \label{yaw rate control: equation}
\end{equation}
where $k_\alpha, k_w>0$ and the measured $w_z$ term acts as an implicit target-yaw-rate feed-forward.

The desired heading $\psi_{\mathrm d}$ is obtained by integrating $\omega_{\mathrm u}$ with angular wrapping. The visibility-conditioned lean $\boldsymbol y^\star$ and $\psi_{\mathrm d}$ define the desired attitude $R_{\mathrm d}\in SO(3)$ through the standard thrust-direction/heading construction. The retained vertical command provides the thrust magnitude ${}^\mathcal{B}T_u$, while a standard geometric $SO(3)$ controller~\cite{lee2010} generates the body torque ${}^\mathcal{B}\boldsymbol{\tau}_u$ required to track $R_{\mathrm d}$. Together, the resulting thrust and torque form the input $\boldsymbol u$ in \eqref{quadrotor dynamics: equation}, thereby closing the visual-control loop.

\section{STABILITY AND VISIBILITY GUARANTEES}
\label{sec:stability}
We now establish the closed-loop guarantees posed in Problem~\ref{prob:landing}. Theorem~\ref{thm:translational} establishes stability and prescribed-performance preservation for the coupled translational visual loop, with Remark~\ref{rem:visibility_stability} extending the result to the visibility-conditioned command. Theorem~\ref{thm:visibility} establishes conditional visibility invariance, while Remark~\ref{rem:yaw_stability} addresses the image-orientation loop.

\begin{theorem}[Coupled Translational Visual Stability]
\label{thm:translational}
Under Assumptions~\ref{asm:visual_regularity} and \ref{asm:uncertainty}, let the initial image-position and image-velocity errors lie strictly inside their prescribed performance envelopes. For any $0<\varphi_{2,\nu}<\varphi_{1,\nu}$, the closed-loop dynamics drive $\boldsymbol\sigma$ exponentially toward $\{\|\boldsymbol\sigma\|\le\vartheta_\nu\}$, where
\begin{equation*}
    \varphi_{1,\nu} \triangleq \underline\beta\, \frac{\min\{2\lambda_{\min}(\Gamma), \lambda_{\min}(\mathcal P)\}}{\max\{1,\underline\beta\lambda_{\max}(\mathcal N^{-1})\}}, ~ \vartheta_\nu \triangleq \sqrt{\frac{3\,\underline\beta\,\lambda_{\max}(\mathcal P)}{\varphi_{1,\nu}-\varphi_{2,\nu}}}\,\tilde d_\nu.
    % \label{translational UUB bound: equation}
\end{equation*}
The adaptive gains $\boldsymbol\kappa$ are also uniformly ultimately bounded. 
Moreover, if $\rho_{\nu,k}(t) > |e_{\nu,z}(t)|\,\bar{s}_k(t)$ for $k\in\{x,y\}$, where $\bar s_k(t) \triangleq |s_{\mathrm d,k}| + \Phi_{kk} \rho_{p,k}(t)$, then $\boldsymbol\zeta_p$ and $\boldsymbol\zeta_\nu$ remain bounded and the prescribed image-position and image-velocity envelopes are preserved.
\end{theorem}

\begin{proof}
Let $\widetilde{\boldsymbol\kappa}=\boldsymbol\kappa-\tilde d_\nu\boldsymbol1$ and consider
\begin{equation}
    V = \frac{1}{2}\boldsymbol\sigma^\top\boldsymbol\sigma + \frac{\underline\beta}{2}\widetilde{\boldsymbol\kappa}^{\top} \mathcal N^{-1}\widetilde{\boldsymbol\kappa}.
    \label{translational Lyapunov: equation}
\end{equation}
The adaptive law preserves $\boldsymbol\kappa\in\mathbb R_{>0}^3$.
Using \eqref{closed loop sliding dynamics: equation}, \eqref{adaptive gain law: equation}, $\beta(t)\ge\underline\beta$, and the row-wise uncertainty bound $|(\mathcal R_\nu\boldsymbol\delta_\nu)_k| \le w_{\nu,k}\tilde d_\nu$, we obtain
\begin{align}
    \dot V \le{}& -\underline\beta\lambda_{\min}(\Gamma)  \|\boldsymbol\sigma\|^2 -\frac{\underline\beta}{2}  \widetilde{\boldsymbol\kappa}^{\top}  \mathcal P\widetilde{\boldsymbol\kappa} + \frac{3}{2} \underline\beta\,\lambda_{\max}(\mathcal P)(\tilde d_\nu)^2 \nonumber\\
    &\le -\varphi_{1,\nu}V + \frac{3}{2} \underline\beta\,\lambda_{\max}(\mathcal P)(\tilde d_\nu)^2,
    \label{translational Vdot bound: equation}
\end{align}
where $\widetilde{\boldsymbol\kappa}^{\top}\mathcal P\boldsymbol\kappa \ge \frac12 \widetilde{\boldsymbol\kappa}^{\top} \mathcal P\widetilde{\boldsymbol\kappa} - \frac12(\tilde d_\nu)^2 \boldsymbol1^\top\mathcal P\boldsymbol1$, and $\boldsymbol1^\top\mathcal P\boldsymbol1 \le 3\lambda_{\max}(\mathcal P)$.
Hence $\dot V\le-\varphi_{2,\nu}V$ whenever
\begin{equation*}
    V \ge \frac{\tfrac{3}{2} \underline\beta\lambda_{\max}(\mathcal P)(\tilde d_\nu)^2}{\varphi_{1,\nu}-\varphi_{2,\nu}},
\end{equation*}
which establishes the stated ultimate bound on $\boldsymbol\sigma$ and uniform ultimate boundedness of $\boldsymbol\kappa$.

For the vertical channel, let $\xi_z \triangleq \int_0^t \zeta_{\nu,z}(\tau)\,d\tau$. Since $\dot{\xi}_z = -\chi_z\xi_z+\sigma_z$, bounded $\sigma_z$ implies bounded $\xi_z$, and hence $\zeta_{\nu,z}=\sigma_z-\chi_z\xi_z$ is bounded.
For each lateral axis $k\in\{x,y\}$,
$\zeta_{\nu,k}=\sigma_k-\chi_k\zeta_{p,k}$, and
\eqref{coupled transformed position dynamics: equation} gives
\begin{equation}
    \dot{\zeta}_{p,k} = -k_p\zeta_{p,k} + \frac{g_{p,k}}{\Phi_{kk}} \left[\rho_{\nu,k}S_{\nu,k} -e_{\nu,z}s_k \right].
    \label{reduced position transformed dynamics: equation}
\end{equation}

On the maximal interval for which the prescribed-position transformation is defined, $|s_k(t)|\leq\bar s_k(t)$. Since $\zeta_{\nu,k}=\sigma_k-\chi_k\zeta_{p,k}$ and $S_{\nu,k}=\tanh(\zeta_{\nu,k}/2)$, bounded $\sigma_k$ implies $S_{\nu,k}\to\mp1$ as $\zeta_{p,k}\to\pm\infty$. Therefore, if $\rho_{\nu,k}(t)>|e_{\nu,z}(t)|\,\bar s_k(t)$ throughout the landing interval, the bracketed term is inward pointing at both extremes. Together with $-k_p\zeta_{p,k}$, this implies bounded $\zeta_{p,k}$. Consequently, $\boldsymbol\zeta_p$ and $\boldsymbol\zeta_\nu$ remain bounded, and the prescribed-performance envelopes are preserved.
\end{proof}

\begin{theorem}[Visibility Forward Invariance]
\label{thm:visibility}
Under Assumption~\ref{asm:visibility_feasibility}, if the measured marker centre is visible initially, i.e., ${}^{\mathcal C}\tilde{\boldsymbol r}_0 \in\mathcal S_{\mathrm{vis}}$, then $\mathcal S_{\mathrm{vis}}$ is forward invariant under the nominal one-step predictor \eqref{predicted marker centre: equation}.
\end{theorem}

\begin{proof}
For ${}^{\mathcal C}\tilde{\boldsymbol r}_j\in\mathcal S_{\mathrm{vis}}$ at discrete prediction step  $j\in\mathbb Z_{\geq0}$, let $\boldsymbol y_j^\star$ denote the solution of \eqref{cbf qp: equation} at that step. Feasibility of \eqref{cbf qp: equation} gives $\left|\left[\Phi^{-1} \boldsymbol c_{\mathrm{next}}(\boldsymbol y_j^\star)\right]_k\right| \le1$ for $k\in\{1,2\}$. Since the nominal predictor \eqref{predicted marker centre: equation} gives ${}^{\mathcal C}\tilde{\boldsymbol r}_{j+1} = \boldsymbol c_{\mathrm{next}}(\boldsymbol y_j^\star)$, we have ${}^{\mathcal C}\tilde{\boldsymbol r}_{j+1}\in\mathcal S_{\mathrm{vis}}$. The result follows by induction.
\end{proof}

\begin{remark}[Virtual-Compass Yaw Stability]
\label{rem:yaw_stability}
Away from the orientation-wrapping discontinuity, the closed-loop image-orientation dynamics satisfy $\ddot e_\alpha+k_w\dot e_\alpha+k_\alpha e_\alpha = \dot d_\alpha + k_w\left(l_{\alpha,1}w_x + l_{\alpha,2}w_y\right)$.
The forcing term is bounded under Assumptions~\ref{asm:orientation_observability}, ~\ref{asm:visual_regularity} and \ref{asm:uncertainty}. Since $s^2+k_ws+k_\alpha$ is Hurwitz for $k_\alpha,k_w>0$, $e_\alpha$ is uniformly ultimately bounded and converges to zero if the forcing term vanishes asymptotically.
\end{remark}

\section{SIMULATION RESULTS}
\label{sec:experiments}

\subsection{Simulation Setup}

VISTA is evaluated in closed-loop simulation using a thirteen-state
X500-class quadrotor model with a downward-facing monocular camera
tracking feature points. The dynamics are integrated at
$\Delta t=0.01$~s subject to actuator constraints. The simulations
include pixel noise and outliers, wind and turbulence, parametric
uncertainty, and centre-of-gravity perturbations. A touchdown is
classified as \emph{precise} when the terminal horizontal
target-relative error is below $0.08$~m and \emph{soft} when the
relative touchdown speed is below $0.20$~m/s; satisfying both defines
a \emph{soft-precise} landing.

For numerical robustness, the simulated visibility QP includes a quadratically penalized nonnegative slack variable. Theorem~\ref{thm:visibility} applies to the nominal zero-slack optimization. Hence, exact forward invariance corresponds to intervals for which the slack remains zero up to solver tolerance. In the successful runs, the slack remains zero or negligible and activates appreciably only after a trajectory enters a failing regime. It therefore does not materially relax the visibility constraint in the reported successful landings.

Robustness to initial position is evaluated from five target-relative UAV configurations,
\begin{equation*}
{}^\mathcal{I}\boldsymbol r_\mathrm{b}(0) = 
\begin{bmatrix}
    {}^\mathcal{I}x_\mathrm{b}(0)\\{}^\mathcal{I}y_\mathrm{b}(0)\\{}^\mathcal{I}z_\mathrm{b}(0)
\end{bmatrix}\in
\left\{
\begin{bmatrix} 0\\0\\5\end{bmatrix},
\begin{bmatrix} 2\\2\\5\end{bmatrix},
\begin{bmatrix} 2\\-2\\5\end{bmatrix},
\begin{bmatrix} 2\\2\\7\end{bmatrix},
\begin{bmatrix} 2\\2\\3\end{bmatrix}
\right\} {\rm m}.   
\end{equation*}
Each initial condition is tested against five target-motion profiles with inertial-frame translational velocities
\[
\begin{aligned}
{}^{\mathcal I}\boldsymbol v_{\mathrm t}^{(1)}
    &= [0,\,0,\,0]^\top,\\
{}^{\mathcal I}\boldsymbol v_{\mathrm t}^{(2)}
    &= [1.1,\,1.1,\,0.1\cos(0.5t)]^\top,\\
{}^{\mathcal I}\boldsymbol v_{\mathrm t}^{(3)}
    &= [0.4\cos(0.8t),\,0.5,\,0]^\top,\\
{}^{\mathcal I}\boldsymbol v_{\mathrm t}^{(4)}
    &= [-0.2\cos(0.5t),\,0.68\cos(0.85t),\,0]^\top,\\
{}^{\mathcal I}\boldsymbol v_{\mathrm t}^{(5)}
    &= [0.24\sin(0.48t),\,0.24\cos(0.48t),
        \,0.1\cos(0.5t)]^\top
        \mathrm{m/s},
\end{aligned}
\]
corresponding to static, linear ship-deck, sinusoidal, Lissajous, and
circular ship-deck motions, respectively. For Cases~2 and~5, the
platform additionally undergoes
$\phi_{\mathrm t}=15^\circ\sin(0.9t)$ and
$\theta_{\mathrm t}=8^\circ\sin(0.6t+\pi/3)$.
Case~2 has $\psi_{\mathrm t}=0$, whereas Case~5 rotates at
$\dot\psi_{\mathrm t}=0.48$~rad/s.


\subsection{Landing Performance and Robustness}
\label{subsec:robustness}

VISTA achieves soft-precise touchdown for all $25$ combinations of
initial condition and target-motion profile without controller
retuning.

\begin{figure}[!t]
\centering
\includegraphics[width=\columnwidth]
{Figures/generated/multi_init/Circular_combined.pdf}
\caption{\textbf{Robustness to five initial conditions for Case~5.} (a) UAV/target trajectories: the shaded region is the $0.08$~m precision region. Touchdown markers denote $\blacktriangle$ for soft-precise touchdown and $\times$ for aborted runs. (b) Camera-frame image-plane trajectories: the gray marker is the desired configuration, while the asymmetric stub encodes image orientation.
The inset magnifies the terminal position and orientation alignment about the desired marker. The legend gives the terminal marker-centre offset $\|\delta{}^{\mathcal C}\hat{\boldsymbol r}\|_2$.}
\label{fig:circular_robustness}
\end{figure}

For the representative rotating Case~5, Fig.~\ref{fig:circular_robustness} shows successful 3-D touchdown from all five initial conditions. In Fig.~\ref{fig:circular_robustness}(b), the marker centres converge toward the desired image location, while the asymmetric stubs align
with the desired gray stub, illustrating the image-orientation regulation analyzed in Remark~\ref{rem:yaw_stability}. The terminal marker-centre offsets are $4.0$--$5.4$~px.

Robustness to target-motion speed is evaluated over the four moving
cases using
$\gamma_v\in\{0.6,0.8,1.0,1.2,1.4\}$, corresponding to a
$\pm40\%$ variation about the nominal motion. All $20$ runs achieve
soft-precise touchdown with unchanged controller parameters; the
worst-case terminal horizontal error and relative touchdown speed are
$4.39$~cm and $0.090$~m/s, respectively.


\subsection{Comparison with Existing Methods}
\label{subsec:comparison}

VISTA is compared with the performance-constrained PBVS (PBVS--PPC) method~\cite{lin2022}, adaptive-observer PBVS (PBVS--AEDO)~\cite{zhang2026}, prescribed-performance IBVS (IBVS--PPC)~\cite{lin2023}, and feed-forward IBVS (FF--IBVS)~\cite{cho2022}. All methods are evaluated from the common initial condition $(2,2,5)$~m under identical target trajectories and disturbance conditions and use the same geometric attitude-control layer.

\begin{figure}[!t]
\centering
\includegraphics[width=\columnwidth]
{Figures/generated/comparison_qual_combined.pdf}
\caption{\textbf{Closed-loop comparison from $\mathrm{IC}_2=(2,2,5)$~m.}
(a) Outcomes over the five target-motion cases, where S denotes soft-precise, H denotes hard or imprecise, and A denotes aborted.
(b) Normalized touchdown energy $\eta_E=(v_f/0.20)^2$: $\log_{10}(\eta_E)<0$ denotes soft touchdown.
(c) UAV landing trajectories for Case~5. Touchdown markers: $\blacktriangle$ soft-precise, $\Diamond$ precise only, $\bigcirc$ soft only, $\bigtriangledown$ neither, and $\times$ aborted.
(d) Normalized marker-centre trajectories: $|{}^{\mathcal C}\tilde r_x|=1$ or $|{}^{\mathcal C}\tilde r_y|=1$ is the FoV boundary.}
\label{fig:controller_comparison}
\end{figure}

Figure~\ref{fig:controller_comparison}(a) shows that VISTA achieves soft-precise touchdown in all five cases. PBVS--PPC, IBVS--PPC, and FF--IBVS abort before reaching the surface in every case, whereas PBVS--AEDO reaches the surface in Cases~1, 2, and~5 but produces hard or imprecise touchdowns.

For Case~5, Fig.~\ref{fig:controller_comparison}(d) shows that VISTA keeps both normalized marker-centre coordinates within the FoV bounds, whereas PBVS--PPC and FF--IBVS reach the camera boundary; PBVS--AEDO remains visible but fails the soft-precise criterion. Using \eqref{relative touchdown energy: equation}, Fig.~\ref{fig:controller_comparison}(b) shows that VISTA remains below the softness threshold in all five cases, while the surface-reaching PBVS--AEDO runs exceed it.

Repeating the comparison over the same $\gamma_v\in\{0.6,0.8,1.0,1.2,1.4\}$ sweep yields no soft-precise touchdown among the $80$ baseline runs. PBVS--AEDO reaches the surface in $8/20$ runs but violates the precision and/or softness requirements, whereas PBVS--PPC, IBVS--PPC, and FF--IBVS do not complete the landing. VISTA achieves soft-precise touchdown in all $20$ corresponding runs without retuning.

\noindent\textit{Discussion.}
These results emphasize the joint landing capability rather than isolated improvements in image alignment, touchdown regulation, or visibility. Because the lateral image-position dynamics are coupled to the vertical image-velocity regulation responsible for soft touchdown, satisfying one objective does not imply satisfaction of the others. VISTA jointly regulates image position, image orientation, and touchdown motion while conditioning the commanded lean for visibility, using unchanged parameters across the tested initial conditions and target-speed variations.

\section{CONCLUSION}
\label{sec:conclusion}

VISTA addresses the intrinsically coupled monocular visual regulation problem in which image-position dynamics depend on the same image-velocity state whose regulation produces soft touchdown, while the attitude corrections required for alignment affect target visibility. Its metric-scale-free architecture jointly regulates image position and depth-normalized relative velocity, yielding lateral alignment and vanishing relative touchdown velocity while preserving orientation and visibility objectives without metric target-depth, metric UAV position/velocity, target-velocity, or absolute-heading estimation.

The closed-loop analysis establishes uniform ultimate boundedness of the coupled visual sliding dynamics and image-orientation error, preservation of the prescribed image-position and image-velocity envelopes under the derived coupling condition, and conditional forward invariance of the visibility set under the nominal one-step predictor. In simulation, VISTA achieved soft-precise touchdown for all $25$ initial-condition/motion combinations and all $20$ runs over the $\pm40\%$ target-speed sweep without retuning. None of the four representative baselines achieved soft-precise touchdown across the corresponding comparisons, underscoring the non-triviality of satisfying alignment, touchdown-severity, orientation, and visibility requirements simultaneously. Future work will address hardware validation under sensing/actuation delays, predictor mismatch, and moving-platform uncertainties.

\bibliographystyle{IEEEtran}
\bibliography{bibliography}

\end{document}